\documentclass[11pt]{amsart}
\usepackage{amsmath,amssymb,amsthm}
\usepackage[margin=1.1in]{geometry}
\usepackage{hyperref}
\usepackage[capitalise,nameinlink]{cleveref}
\emergencystretch=2em

\newtheorem{theorem}{Theorem}
\newtheorem{lemma}[theorem]{Lemma}
\theoremstyle{remark}
\newtheorem{remark}[theorem]{Remark}

\title{Which chord of an even cycle preserves its spectral gap}
\author{Yitzchak Shmalo}
\date{}

\begin{document}

\begin{abstract}
Add one edge of conductance $w>0$ to the cycle $C_n$, $n=2m\ge4$, and run the random walk on the resulting weighted
graph. We show that the spectral gap of the walk is largest exactly when the new edge joins two antipodal vertices,
and that in this case the gap equals $1-\cos(\pi/m)$, the gap of the cycle itself, for every $w$. Every other chord
strictly lowers the gap, with an explicit bound. If instead the walk is Metropolized so that its stationary law stays
uniform, every chord, antipodal or not, strictly lowers the gap. These are one-edge cases of questions of Diaconis
about which added edges speed up a random walk most.
\end{abstract}

\maketitle

\section{Introduction}

Diaconis~\cite[Section~5, Question~4]{Diaconis} asks, for a connected regular graph, which perfect matching should
be added to make the random walk mix fastest, and, for a general reversible chain, which weighted edge should be added,
with the new chain Metropolized to keep its stationary law, to improve mixing or the spectral gap the most. For the
cycle with a random perfect matching, the mixing behaviour is now well understood; see Hermon, Sly and
Sousi~\cite{HSS} and the references there. It is also known that adding an edge can lower the spectral gap of the
random walk; Eldan, R\'acz and Schramm~\cite{ERS} show that this happens with substantial probability in random graphs.
For the combinatorial Laplacian an added edge never lowers the algebraic connectivity, and the choice of edges
that raise it most has been studied by Ghosh and Boyd~\cite{GB} and by Kim~\cite{Kim}. For a single weighted chord
on a weighted cycle, Deng et al.~\cite{DCY} give exact formulas for the effective resistances
and the Kirchhoff index after the chord is added, and show that under bounded conductances near-antipodal chords that
balance the resistances of the two arcs are near-optimal for the algebraic connectivity.

Here we treat the smallest case completely: a single chord on an even cycle. Write $n=2m\ge4$, label the vertices
$\mathbb Z/n$, give every cycle edge conductance $1$, and add one chord $\{0,k\}$, $2\le k\le m$, of conductance $w>0$
(by rotation and reflection every chord has this form). Let $d_i$ be the weighted degrees, $A$ the weighted adjacency
matrix, $P=D^{-1}A$ the random walk, and
\[
\gamma=1-\lambda_2(P)=\inf_f\frac{\mathcal E(f)}{\operatorname{Var}_D(f)},\qquad
\mathcal E(f)=\sum_{\{i,j\}}c_{ij}(f_i-f_j)^2,\qquad \operatorname{Var}_D(f)=\min_{c\in\mathbb R}\sum_id_i(f_i-c)^2,
\]
where the sum runs over edges, each counted once, and the infimum over nonconstant $f$. The cycle alone has
$\gamma(C_n)=1-\cos(2\pi/n)$. Throughout, $a=1-\cos(\pi/m)=1-\cos(2\pi/n)$.

\begin{theorem}\label{thm:walk}
For every $w>0$, $\gamma(C_{2m}+w\{0,m\})=a$. For $2\le k<m$,
\begin{equation}\label{eq:strict}
\gamma(C_n+w\{0,k\})\le\frac{a}{1+\frac{2w}{n+w}\cos^2(\pi k/n)}<a .
\end{equation}
In particular the antipodal chords, and only they, maximise the spectral gap among all single chords of a given
weight, and no chord raises the gap of the cycle.
\end{theorem}

For the second question, take as proposal the walk $P$ above and as target the uniform law. The Metropolis--Hastings
kernel is $K(i,j)=P(i,j)\min\{1,P(j,i)/P(i,j)\}=c_{ij}/\max(d_i,d_j)$ for $i\ne j$, with the diagonal chosen so that
rows sum to one. It is symmetric, so the uniform law is stationary.

\begin{theorem}\label{thm:mh}
For every chord $\{0,k\}$, $2\le k\le m$, and every $w>0$, with $b=w/(2+w)$, $\theta=2\pi/n$ and $\alpha=\pi k/n$,
\[
1-\lambda_2(K)\le a-\frac{4b}{n}\bigl(a^2\cos^2\alpha+\sin^2\theta\,\sin^2\alpha\bigr)<a .
\]
So the Metropolized chain is strictly slower, in the sense of the spectral gap, than the walk on the cycle without the
chord, whichever chord is added.
\end{theorem}

The antipodal case is the point of \cref{thm:walk}: the upper bound for the other chords comes from one test function,
but the equality needs a lower bound on the whole Rayleigh quotient. The question for full matchings remains open.

\section{The random walk}

\begin{lemma}\label{lem:path}
If $g_0,\dots,g_m$ are real with $g_0=g_m=0$, then $\sum_{i=0}^{m-1}(g_{i+1}-g_i)^2\ge2a\sum_{i=1}^{m-1}g_i^2$.
\end{lemma}

\begin{proof}
The left side is $g^{\top}Tg$ for the $(m-1)\times(m-1)$ tridiagonal matrix $T$ with $2$ on the diagonal and $-1$ next
to it. Its eigenvectors are $(\sin(\pi ji/m))_i$, with eigenvalues $2-2\cos(\pi j/m)$, $1\le j\le m-1$, the smallest
being $2a$.
\end{proof}

\begin{proof}[Proof of \cref{thm:walk}, antipodal chord]
Let $k=m$, so that $d_0=d_m=2+w$ and $d_i=2$ otherwise. The half-turn $(Tf)(i)=f(i+m)$ preserves the graph and the
degrees, so it preserves $\mathcal E$ and the $D$-inner product. Write $f=f_++f_-$ with $Tf_\pm=\pm f_\pm$. The two
parts are orthogonal for both forms, and $\sum_id_if_-(i)=0$, so
\[
\mathcal E(f)=\mathcal E(f_+)+\mathcal E(f_-),\qquad \operatorname{Var}_D(f)=\operatorname{Var}_D(f_+)+\sum_id_if_-(i)^2 .
\]

Even part. Put $t=f_+(0)=f_+(m)$ and $g=f_+-t$. The chord contributes nothing to $\mathcal E(f_+)$, and $g$ vanishes at
both ends of each of the two arcs from $0$ to $m$. \Cref{lem:path} on both arcs gives
$\mathcal E(f_+)\ge2a\sum_ig_i^2$. Since $g$ vanishes where $d_i\ne2$,
$\operatorname{Var}_D(f_+)\le\sum_id_ig_i^2=2\sum_ig_i^2$, so $\mathcal E(f_+)\ge a\operatorname{Var}_D(f_+)$.

Odd part. Put $s=f_-(0)=-f_-(m)$. The function $f_-$ has $\sum_if_-(i)=0$, so the Poincar\'e inequality of the cycle,
whose Laplacian has second eigenvalue $2-2\cos(2\pi/n)=2a$, gives
$\mathcal E_{\mathrm{cyc}}(f_-)\ge2a\sum_if_-(i)^2$. The chord adds $w(2s)^2$. Since $a\le1$,
\[
\mathcal E(f_-)\ge2a\sum_if_-(i)^2+4ws^2\ge a\Bigl(2\sum_if_-(i)^2+2ws^2\Bigr)=a\sum_id_if_-(i)^2 .
\]
Adding the two parts, $\mathcal E(f)\ge a\operatorname{Var}_D(f)$, so $\gamma\ge a$.

For the reverse inequality take $f(i)=\sin(\pi i/m)$. It vanishes at $0$ and $m$, so the chord does not affect
$(D-A)f$ and $(D-A)f=2af=aDf$. Also $\sum_id_if(i)=2\sum_if(i)=0$. Hence $f$ is an eigenvector of $P$ with eigenvalue
$1-a$ that is orthogonal to the constants, and $\gamma\le a$.
\end{proof}

\begin{proof}[Proof of \cref{thm:walk}, other chords]
Let $2\le k<m$, $\theta=2\pi/n$, $t=\cos(\pi k/n)>0$ and $f(i)=\cos(\theta(i-k/2))$. Then $f(0)=f(k)=t$, so the
chord contributes no energy, and $f$ is an eigenvector of the cycle Laplacian with eigenvalue $2a$, with
$\sum_if(i)=0$ and $\sum_if(i)^2=n/2$. Hence $\mathcal E(f)=na$. With $\sum_id_if(i)^2=n+2wt^2$,
$\sum_id_if(i)=2wt$ and $\sum_id_i=2n+2w$,
\[
\operatorname{Var}_D(f)=n+2wt^2-\frac{(2wt)^2}{2n+2w}=n\Bigl(1+\frac{2wt^2}{n+w}\Bigr),
\]
and $\gamma\le\mathcal E(f)/\operatorname{Var}_D(f)$ is~\eqref{eq:strict}.
\end{proof}

\begin{remark}
The same test function gives the bound~\eqref{eq:strict} for every chord of an odd cycle, with $a=1-\cos(2\pi/n)$, since
the argument uses only $\cos(\pi k/n)>0$, which holds for $2\le k\le(n-1)/2$.
\end{remark}

\section{The Metropolized walk}

\begin{proof}[Proof of \cref{thm:mh}]
The chord endpoints have degree $2+w$ and every other vertex degree $2$. So $K$ has conductance $1/(2+w)$ on the four
cycle edges at the chord endpoints (call this set $S$), $1/2$ on the other cycle edges, and $w/(2+w)=b$ on the chord.
Since $1/(2+w)=\tfrac12-\tfrac b2$, for every real $f$,
\[
f^{\top}(I-K)f=\tfrac12\mathcal E_{\mathrm{cyc}}(f)-\tfrac b2\sum_{\{i,j\}\in S}(f_i-f_j)^2+b(f_0-f_k)^2 .
\]
The four edges of $S$ are distinct because $2\le k\le n-2$. Take again $f(i)=\cos(\theta i-\alpha)$, which has mean
zero, $\sum_if(i)^2=n/2$, $\mathcal E_{\mathrm{cyc}}(f)=na$ and $f(0)=f(k)=\cos\alpha$. Its values at the neighbours
of the endpoints are $f(\pm1)=f(k\mp1)$, equal to $\cos(\theta\mp\alpha)$, and a short computation gives
\[
\sum_{\{i,j\}\in S}(f_i-f_j)^2=2\bigl[(a\cos\alpha-\sin\theta\sin\alpha)^2+(a\cos\alpha+\sin\theta\sin\alpha)^2\bigr]
=4\bigl(a^2\cos^2\alpha+\sin^2\theta\sin^2\alpha\bigr).
\]
Dividing $f^{\top}(I-K)f$ by $\sum_if(i)^2=n/2$ gives the bound. The bracket is positive because $a>0$ and
$\sin\theta>0$ for $n\ge4$, so the bound is below $a$, the gap of the walk on $C_n$.
\end{proof}

\begin{remark}
\Cref{thm:mh} fixes one natural proposal, the walk on the augmented graph. Diaconis's question allows any way of
Metropolizing, and among the slower kernels of \cref{thm:mh} we do not determine the best chord.
\end{remark}

\begin{remark}
Exact computations agree with both theorems. For even $n\le16$, every chord and $w\in\{1/2,1,5\}$, the characteristic
polynomial of the pencil $(D-A)-xD$ was computed in rational arithmetic in two independent ways, by direct
determinants and by eliminating the two path interiors, and the position of its second root was certified by Sturm
sequences: $21$ equalities and $63$ strict inequalities, as predicted. Floating-point checks for even $n\le40$, five
weights from $0.01$ to $1000$, and every chord, and for odd $n\le29$, agree with~\eqref{eq:strict} and
\cref{thm:mh}. The code is in the repository listed below.
\end{remark}

\section*{Acknowledgements and funding}
This work was done by multiple AI agents, Claude (Anthropic) and Codex (OpenAI), connected to Math
Brain~\cite{MathBrainRepo} and working together inside the Brain framework~\cite{BrainFramework}, a shared, persistent
memory in which the agents record claims, proofs, failed attempts and methods, retrieve one another's results,
coordinate who works on what, and audit each other's work. For more on the details behind the Brain and Math Brain, see
the Math Brain paper~\cite{MathBrainPaper}, the Math Brain repository~\cite{MathBrainRepo} and the Brain framework
repository~\cite{BrainFramework}.

The research presented in this paper was supported by the European Research Council (ERC) under the European Union's
Horizon Europe research and innovation programme (grant agreement No.~101041711), by the Simons Foundation as part of
the Collaboration on the Mathematical and Scientific Foundations of Deep Learning, by Heights Labs, by Convex Nexus
Capital, by the Israel Science Foundation (grant number 2258/19), and by the Israel Science Foundation (ISF Grant
4101/25). Computational resources were provided in part by Google Cloud through its education credits program (credit
530550103).

\section*{Declaration of competing interest}
The author holds equity in Heights Labs and in Convex Nexus Capital and has a management role at Convex Nexus Capital.
These entities are included in the funding declaration above.

\section*{Declaration of generative AI and AI-assisted technologies}
The paper comes out of Math Brain, an ongoing project on the use of AI systems for research mathematics, and its
results were developed with the Brain~\cite{BrainFramework} and Math Brain~\cite{MathBrainRepo}. In September and
November 2026 Claude agents (Anthropic) and Codex agents running OpenAI's GPT-6.1, connected to the Brain, studied
Diaconis's question. A Codex agent proved both theorems and wrote the exact checks; Claude agents audited the proofs
line by line and checked two algebraic steps of the proof of \cref{thm:mh} in Lean, and Claude Opus 5.5 (Anthropic) re-derived the proofs,
checked them numerically, and prepared the manuscript. These checks are not peer review. The author directed the work
and is responsible for the content of the paper and for decisions concerning its dissemination.

\section*{Code and data availability}
Code and exact computations are at

\url{https://github.com/yspennstate/one-chord-even-cycle-spectral-gap}.

The paper is part of the Math Brain project, whose public repository is
\url{https://github.com/yspennstate/math-brain-public}~\cite{MathBrainRepo}.

\begin{thebibliography}{99}
\bibitem{DCY} J.~Deng, Liu Chang, Q.~Yang, Single-chord augmentation of weighted cycles for algebraic connectivity
and network coherence, preprint, 2026, arXiv:2605.24479.
\bibitem{Diaconis}P.~Diaconis, Some things we've learned (about Markov chain Monte Carlo), Bernoulli 19 (2013),
1294--1305, arXiv:1309.7754.
\bibitem{ERS} R.~Eldan, M.~Z.~R\'acz, T.~Schramm, Braess's paradox for the spectral gap in random graphs and
delocalization of eigenvectors, Random Structures Algorithms 50 (2017), 584--611.
\bibitem{GB} A.~Ghosh, S.~Boyd, Growing well-connected graphs, in: Proc. 45th IEEE Conference on Decision and
Control, 2006, 6605--6611.
\bibitem{HSS} J.~Hermon, A.~Sly, P.~Sousi, Universality of cutoff for graphs with an added random matching, Ann.
Probab. 50 (2022), 203--240, arXiv:2008.08564.
\bibitem{Kim} Y.~Kim, Bisection algorithm of increasing algebraic connectivity by adding an edge, IEEE Trans.
Automat. Control 55 (2010), 170--174.
\bibitem{MathBrainRepo} Y.~Shmalo, \emph{Math Brain: public snapshot, code and tests}, 2026,
\url{https://github.com/yspennstate/math-brain-public}.
\bibitem{MathBrainPaper} Y.~Shmalo, \emph{Math Brain: shared, auditable memory for multi-agent mathematical research},
October 2026, \url{https://github.com/yspennstate/math-brain-public/tree/main/papers/math-brain-2026-10-24}.
\bibitem{BrainFramework} Y.~Shmalo, \emph{Brain Framework}, public software repository, 2026,
\url{https://github.com/yspennstate/brain-framework-public}.
\end{thebibliography}

\end{document}
